\chapter*{ضمیمهٔ ۳: واژه‌نامهٔ تخصصی}
\addcontentsline{toc}{chapter}{ضمیمهٔ ۳: واژه‌نامه}

\begin{center}
\textit{تعریف مختصر اصطلاحات کلیدی کتاب}
\end{center}

\vspace{5mm}



\textbf{آنارشیسم} (\en{Anarchism}): فلسفهٔ سیاسی که خواستار لغو دولت و هر نوع سلسله‌مراتب اجباری است.

\textbf{آنارکو-کاپیتالیسم} (\en{Anarcho-Capitalism}): شکلی از آنارشیسم که خواستار لغو دولت و جایگزینی آن با بازار آزاد است.

\textbf{آنارکو-سندیکالیسم} (\en{Anarcho-Syndicalism}): آنارشیسمی که بر اتحادیه‌های کارگری به عنوان ابزار انقلاب تأکید دارد.

\textbf{اتوریتارینیسم} (\en{Authoritarianism}): نظام سیاسی با تمرکز قدرت و محدودیت آزادی‌ها، اما بدون کنترل کامل بر جامعه (برخلاف توتالیتاریسم).

\textbf{اصل عدم تجاوز} (\en{Non-Aggression Principle/NAP}): اصل لیبرتارین که می‌گوید استفادهٔ ابتدایی از زور علیه دیگران نادرست است.

\textbf{اقتصاد بازار اجتماعی} (\en{Social Market Economy}): مدل آلمانی که ترکیبی از بازار آزاد و حمایت اجتماعی است.

\textbf{اقتصاد مختلط} (\en{Mixed Economy}): نظام اقتصادی که ترکیبی از عناصر سرمایه‌دارانه و سوسیالیستی دارد.

\textbf{انترناسیونالیسم} (\en{Internationalism}): باور به همکاری و همبستگی فراملی، معمولاً در چپ.

\textbf{برابری‌خواهی} (\en{Egalitarianism}): باور به برابری اساسی انسان‌ها و لزوم کاهش نابرابری.

\textbf{بورژوازی} (\en{Bourgeoisie}): در مارکسیسم، طبقهٔ مالک ابزار تولید؛ سرمایه‌داران.

\textbf{پرولتاریا} (\en{Proletariat}): در مارکسیسم، طبقهٔ کارگر که نیروی کار خود را می‌فروشد.

\textbf{پوپولیسم} (\en{Populism}): سبک سیاسی که «مردم» را در برابر «نخبگان» قرار می‌دهد.

\textbf{تکثرگرایی} (\en{Pluralism}): باور به ارزش تنوع گروه‌ها و دیدگاه‌ها در جامعه.

\textbf{توتالیتاریسم} (\en{Totalitarianism}): نظام سیاسی که کنترل کامل بر همهٔ جنبه‌های زندگی دارد.

\textbf{جامعهٔ مدنی} (\en{Civil Society}): حوزهٔ نهادها و انجمن‌های داوطلبانه بین خانواده و دولت.

\textbf{جمهوری‌خواهی} (\en{Republicanism}): سنت فکری که بر مشارکت شهروندی، آزادی از سلطه، و خیر عمومی تأکید دارد.

\textbf{چپ} (\en{Left}): طیف سیاسی که تاریخاً بر برابری، تغییر، و عدالت اجتماعی تأکید دارد.

\textbf{حقوق طبیعی} (\en{Natural Rights}): حقوقی که انسان‌ها به صرف انسان بودن دارند، مستقل از قانون موضوعه.

\textbf{دموکراسی} (\en{Democracy}): نظام حکومتی که در آن قدرت از مردم ناشی می‌شود.

\textbf{دموکراسی اجتماعی} (\en{Social Democracy}): ایدئولوژی که خواستار اصلاح سرمایه‌داری از طریق دولت رفاه و دموکراسی است.

\textbf{دموکراسی مسیحی} (\en{Christian Democracy}): جریان سیاسی مبتنی بر آموزهٔ اجتماعی مسیحی، معمولاً میانه‌روی.

\textbf{دولت رفاه} (\en{Welfare State}): دولتی که مسئولیت تأمین رفاه اجتماعی شهروندان را می‌پذیرد.

\textbf{راست} (\en{Right}): طیف سیاسی که تاریخاً بر سنت، سلسله‌مراتب، و نظم تأکید دارد.

\textbf{رفرمیسم} (\en{Reformism}): رویکردی که خواستار تغییر تدریجی از طریق اصلاحات است، نه انقلاب.

\textbf{سرمایه‌داری} (\en{Capitalism}): نظام اقتصادی مبتنی بر مالکیت خصوصی ابزار تولید و بازار آزاد.

\textbf{سکولاریسم} (\en{Secularism}): جدایی دین از سیاست؛ بی‌طرفی دولت نسبت به دین.

\textbf{سوسیالیسم} (\en{Socialism}): ایدئولوژی که خواستار مالکیت اجتماعی بر ابزار تولید است.

\textbf{طبقه} (\en{Class}): در مارکسیسم، گروه‌های اجتماعی که بر اساس رابطه با ابزار تولید تعریف می‌شوند.

\textbf{فاشیسم} (\en{Fascism}): ایدئولوژی افراطی قرن بیستم با ویژگی‌های ناسیونالیسم افراطی، اقتدارگرایی، و ضد لیبرالیسم.

\textbf{فدرالیسم} (\en{Federalism}): نظام سیاسی که قدرت بین دولت مرکزی و واحدهای محلی تقسیم می‌شود.

\textbf{فردگرایی} (\en{Individualism}): باور به اولویت فرد بر جامعه یا دولت.

\textbf{کمونیسم} (\en{Communism}): ایدئولوژی که خواستار جامعهٔ بی‌طبقه و لغو مالکیت خصوصی است.

\textbf{کمونیتارینیسم} (\en{Communitarianism}): فلسفه‌ای که بر اهمیت جامعه و پیوندهای اجتماعی تأکید دارد.

\textbf{کورپوراتیسم} (\en{Corporatism}): سازماندهی اقتصاد بر اساس گروه‌های صنفی با نظارت دولت.

\textbf{لیبرالیسم} (\en{Liberalism}): ایدئولوژی که بر آزادی فردی، حقوق، و حکومت قانون تأکید دارد.

\textbf{لیبرتارینیسم} (\en{Libertarianism}): فلسفه‌ای که بر آزادی فردی حداکثری و دولت حداقلی تأکید دارد.

\textbf{مارکسیسم} (\en{Marxism}): نظریهٔ اجتماعی-اقتصادی-سیاسی مبتنی بر آثار مارکس و انگلس.

\textbf{محافظه‌کاری} (\en{Conservatism}): ایدئولوژی که بر حفظ سنت، نهادها، و تغییر تدریجی تأکید دارد.

\textbf{مینارشیسم} (\en{Minarchism}): لیبرتارینیسمی که دولت حداقلی (فقط امنیت و دادگاه) را می‌پذیرد.

\textbf{ناسیونالیسم} (\en{Nationalism}): ایدئولوژی که بر هویت و حاکمیت ملت تأکید دارد.

\textbf{نئولیبرالیسم} (\en{Neoliberalism}): احیای لیبرالیسم اقتصادی از دهۀ ۱۹۷۰؛ بازار آزاد و کاهش دولت.

\textbf{نئوکنسرواتیسم} (\en{Neoconservatism}): جریان محافظه‌کار آمریکایی با تأکید بر سیاست خارجی مداخله‌جویانه.

\textbf{نئومارکسیسم} (\en{Neo-Marxism}): بازخوانی‌های قرن بیستمی مارکسیسم (مکتب فرانکفورت و دیگران).

\textbf{هژمونی} (\en{Hegemony}): در گرامشی، سلطهٔ فرهنگی-ایدئولوژیک طبقهٔ حاکم.

\textbf{یوروسپتیسیزم} (\en{Euroscepticism}): مخالفت یا شک نسبت به اتحادیهٔ اروپا و یکپارچگی اروپایی.


